problem:  
Let \((X,d,\mathfrak{m})\) be a compact metric measure space satisfying the curvature-dimension condition \(\mathrm{CD}(K,N)\) for some \(K \in \mathbb{R}\), \(1 < N < \infty\).  
For each pair \(\mu_0, \mu_1 \in \mathcal{P}_2(X)\) absolutely continuous with respect to \(\mathfrak{m}\), let \((\mu_t^\varepsilon)_{t \in [0,1]}\) denote the **entropic interpolation** (marginal flow of the Schrödinger problem) with parameter \(\varepsilon > 0\), and let \(\mu_t\) denote the limiting **Wasserstein geodesic interpolation** as \(\varepsilon \to 0\).

Assume that for some sequence \(\varepsilon_j \to 0\), the family of “renormalized entropic actions”
\[
\mathcal{A}_{\varepsilon_j}(\mu_0,\mu_1) \ :=\ \frac{1}{\varepsilon_j}\int_0^1\Big(\mathrm{Ent}_{\mathfrak{m}}(\mu_t^{\varepsilon_j}) - \mathrm{Ent}_{\mathfrak{m}}(\mu_t)\Big)\,dt
\]
converges for all absolutely continuous \(\mu_0,\mu_1\) to a finite, **quadratic functional** \(Q(\mu_0,\mu_1)\) on \(\mathcal{P}_2(X)\).

**Question:**  
Does the existence of such a finite quadratic limit functional \(Q\) for all pairs \(\mu_0,\mu_1\) necessarily imply that \((X,d,\mathfrak{m})\) is an \(\mathrm{RCD}(K,N)\) space, and moreover that \(Q\) coincides (up to normalization) with the integrated Bochner–type “curvature term” \(\int_0^1\!\!\int_X \langle \mathrm{Ric}_X(\nabla \varphi_t,\nabla \varphi_t)\rangle\,d\mathfrak{m}\,dt\) in Gigli’s second-order calculus?  
In other words: **can the asymptotic defect between Schrödinger and Wasserstein transports detect the existence and the tensorial form of the Ricci curvature in an otherwise synthetic \(\mathrm{CD}(K,N)\) space?**

Why is it a "good" problem:  
This refined formulation focuses the core question into a logically coherent bridge between **entropic optimal transport asymptotics** and **tensorial curvature structures** in non-smooth metric measure spaces.  
   
1. **Definite logical structure:**  
   The problem asks whether a purely analytic signature—a second-order asymptotic of entropy-regularized transport—characterizes precisely the \(RCD(K,N)\) class among CD spaces, thus providing a measurable analytic test of infinitesimal Riemannianity.  
   
2. **Conceptual unification:**  
   It blends three deep currents: the Schrödinger problem from stochastic analysis, the information-geometry viewpoint on entropy, and the Bochner inequality from geometric analysis. Establishing equivalence or detecting curvature tensoriality from entropic asymptotics would represent a new synthesis.  
   
3. **Mathematical depth with simplicity:**  
   The question is stated in a few lines but cuts to the foundational nature of curvature in nonsmooth metric geometry—whether the “heat-fluctuation geometry” of the Schrödinger bridge inherently contains the same second-order curvature information as the Bochner inequality.  
   
4. **Potential for breakthrough:**  
   A positive result would create a new analytic characterization of \(RCD(K,N)\) spaces through transport asymptotics, introducing entirely new methods to probe curvature via entropic action functionals.  
   
Hence, this problem is “good” because it is concise, conceptually central, and aims to connect entropic regularization theory with the tensor-level manifestation of synthetic Ricci curvature, a connection not established by any existing theorem.